This chapter presents the design considerations that guided the development of the implemented IoT- and AI-driven fuel monitoring and anomaly detection system. Whereas Chapter~\ref{ch:theorycon} established the theoretical foundations of IoT architectures, cyber-physical systems, signal processing, and context-aware anomaly detection, the present chapter translates those theories into practical engineering decisions that were directly implemented in Chapter~\ref{ch:method}. 

These design considerations serve as the structural “columns’’ of the system: grounded in accepted standards, ethical guidelines, and engineering best practices, but tailored to support the procedural steps carried out in Chapter~\ref{ch:method}. The discussion begins with a comprehensive account of the technical standards that ensure reliability, safety, interoperability, and data security, followed by software and hardware design decisions. Ethical considerations—including privacy, responsible AI, and driver safety—are also presented.

\section{Standards}
\label{sec:standards}

\subsection{Adherence to Technical Standards}

The prototype was designed and implemented in accordance with the technical standards listed in this section. Each standard was applied as a design constraint that narrowed the space of acceptable hardware components, signal-handling procedures, transport mechanisms, and software structures, and each therefore produced an identifiable decision that is visible in the implementation described in Chapter~\ref{ch:method}. Adherence was established at the design level, which is the level at which the influence of each standard is documented here; formal third-party certification of the assembled device was outside the scope of an undergraduate prototype.

Table~\ref{tab:standards_mapping} maps every standard and regulatory instrument observed in this study to the specific design decision it produced and to the section of Chapter~\ref{ch:method} in which that decision is implemented. The discussion that follows explains the reasoning behind the mappings that are not self-evident from the table.

\FloatBarrier
{\scriptsize
\setlength{\tabcolsep}{3pt}
\renewcommand{\arraystretch}{1.15}
\begin{longtable}{|p{0.17\textwidth}|p{0.15\textwidth}|p{0.45\textwidth}|p{0.14\textwidth}|}
\caption{Mapping of Cited Standards to Implemented Design Decisions}
\label{tab:standards_mapping} \\
\hline
\textbf{Standard / Instrument} & \textbf{Domain Governed} & \textbf{Design Decision Directly Produced} & \textbf{Implemented In} \\
\hline
\endfirsthead

\hline
\textbf{Standard / Instrument} & \textbf{Domain Governed} & \textbf{Design Decision Directly Produced} & \textbf{Implemented In} \\
\hline
\endhead

\stdlink{iso15765_4}{ISO~15765-4}, \stdlink{iso15031_5}{ISO~15031-5}, \stdlink{sae_j1349}{SAE~J1349} & Vehicle data acquisition & Fixed the acquisition path as diagnostics over CAN at the OBD-II port, addressed through an MCP2515 controller and transceiver, with no tank or sender-loom modification. The standardised Mode~01 PID set exposed over this transport defined the contextual feature budget available at zero integration cost -- speed, RPM, engine load, throttle, MAF, coolant temperature, and runtime -- which became the backbone of the 33-feature Model~C vector. The ISO~15765-4 minimum inter-request interval and the SAE~J1349 idle-combustion bound act as opposing constraints that jointly fix the sampling interval at $\Delta s = 1.5$~s. & \seclink{sec:hardware}{Secs.~5.3, 5.4.1--5.4.4} \\
\hline
\stdlink{ieee1451}{IEEE~1451}, NMEA~0183 & Peripheral interfacing & Enforced a uniform digital interface discipline: every sensing element reaches the MCU over a standard bus with no analogue conditioning in the signal path, allowing independent substitution of peripherals. GNSS output is parsed as ASCII sentences over UART rather than a vendor binary protocol, which permits the mobile-application GPS fallback to be fused with the embedded receiver at the backend under a common schema and a \texttt{gps\_source} discriminator. & \seclink{sec:gps_pipeline}{Secs.~5.3, 5.7.1, 5.7.4} \\
\hline
\stdlink{iec61131}{IEC~61131-2}, \stdlink{iso7637}{ISO~7637-2}, \stdlink{cecert}{CE}/\stdlink{fcc}{FCC} marking & Hardware robustness and conformity & Established the immunity and environmental classes the in-cab node was designed to meet, producing a deterministic scan-style main loop with watchdog-backed recovery and a hardware modem power-control line. Required a regulated supply behind the OBD-II 12~V line, and, because transient-induced reset is an anticipated event, required the offline buffer to persist in SPIFFS flash rather than RAM. Constrained procurement to CE/FCC-marked module variants and bounded the LoRa configuration to 433~MHz, SF7, 125~kHz, 17~dBm. & \seclink{sec:hardware}{Secs.~5.3, 5.6.13.1, 5.6.15} \\
\hline
\stdlink{ntc2025legacy_network_phaseout}{NTC~MC No.~003-09-2025} & Wide-area link selection & Determined the modem choice. Because the circular provides for phased 2G and 3G decommissioning, a GSM-only baseline was rejected as a terminal design and the LTE Cat-1 A7670E adopted, placing LTE as the primary path with 2G retained as a transitional fallback on the same device. & \seclink{sec:communication}{Secs.~4.2.4, 5.6.1} \\
\hline
\stdlink{rfc791}{IETF~TCP/IP}, \stdlink{ieee80211}{IEEE~802.11}, TLS/SSL with \stdlink{iso27001}{ISO/IEC~27001} & Transport and communication security & Fixed the transport contract for all wide-area channels as IP over TCP, allowing LTE, GSM, and WiFi to share a single HTTP client and reducing channel selection to a connectivity predicate. The 802.11 four-way handshake defined the WiFi detection predicate and its measured probe cost. Mandated HTTPS for all telemetry with credentials held in device-side configuration; the resulting TLS negotiation is a documented and accepted contributor to the 23--24~s median LTE latency. & \seclink{sec:communication}{Secs.~5.6.2--5.6.6, 5.6.14} \\
\hline
IETF RFC~5905 (NTP), FAIR principles & Time base and data provenance & Established an NTP-disciplined device clock giving a common epoch-millisecond time base across fuel, GPS, speed, and engine channels, and produced the telemetry schema: stable record identifier, monotonic sequence number, dual device-side and backend-side timestamps, and channel provenance. Together these make drift measurable, offline-replayed records distinguishable and correctly re-orderable, and the evaluation corpus independently reproducible. & \seclink{sec:fuel_preprocessing}{Secs.~5.4.6, 5.6.10, 5.9} \\
\hline
\stdlink{iso42010}{ISO/IEC/IEEE~42010} & Software architecture & Produced the five-layer decomposition (acquisition, processing, communication, cloud integration, presentation) with explicit interface contracts between layers, which permitted machine-learning inference to be relocated from the ESP32 to the backend without altering the acquisition or transport layers. & \seclink{sec:implement}{Secs.~4.2.2, 5.2} \\
\hline
ISO/IEC~25010, \stdlink{iso29119}{ISO/IEC~29119} & Quality criteria and test design & Supplied the quality characteristics reported in Chapter~6: functional correctness (fuel MAPE), performance efficiency (latency), reliability (packet delivery and offline recovery); and determined the test architecture: controlled anomaly injection with known ground truth, leave-one-trip-out cross-validation to prevent within-trip leakage, and an out-of-distribution clean-data condition bounding the false-positive rate. & \seclink{sec:evaluation}{Secs.~5.10, 5.11} \\
\hline
\stdlink{iso27018}{ISO/IEC~27018}, GDPR & Personal data protection & Produced schema and access decisions rather than policy statements: driver identifiers stored as opaque references, row-level security enforced at the database layer, role-scoped dashboard views, and aggregate retention of location and behavioural data for reporting. & \seclink{sec:backend}{Sec.~5.8} \\
\hline
\stdlink{iso26262}{ISO~26262}, \stdlink{iso39001}{ISO~39001}, \stdlink{iso45001}{ISO~45001} & Safety and system authority & Established the limits of system authority: the prototype is strictly read-only on the vehicle bus, issues no actuation commands, and presents all driver-facing output as advisory. Produced the rest-alert subsystem as a cumulative-exposure mechanism keyed to continuous driving time and distance with acknowledgement logged as an event, and governed the physical installation as a non-invasive OBD-II connection in an enclosure that does not obstruct the driver's controls or field of view. & \seclink{sec:hardware}{Secs.~4.1.2, 5.3.1, 5.4.1, 5.8.7} \\
\hline
\stdlink{ieeep7003}{IEEE~P7003} & AI accountability & Produced the requirement that anomaly alerts carry the contextual evidence that triggered them -- fuel delta, speed, engine state, trip phase -- rather than an opaque score, and that the contextual gate suppressing alerts during refuelling and post-refuel quiescence be documented so that a flagged driver can be shown why the alert fired. & \seclink{sec:evaluation}{Secs.~5.10, 5.10.11} \\
\hline
\end{longtable}
}
\normalsize
\FloatBarrier

\subsubsection*{Discussion}

The groupings in Table~\ref{tab:standards_mapping} are elaborated below where the reasoning is not self-evident from the table: where a standard eliminated an alternative, where adherence carried a cost the design accepted, or where a standard was evaluated and deferred.

\paragraph{Sensors and hardware.} The partner operator's requirement that no vehicle be physically modified, consistent with ISO~45001~\cite{iso45001}, made a diagnostic-port interface the acquisition path for the system and removed in-tank capacitive and float sensors from consideration before electrical design began. ISO~15765-4~\cite{iso15765_4} then fixed the transport layer and, through the standardised Mode~01 PID set it carries, determined the contextual feature budget available without additional instrumentation. Where the standardised PID set proved insufficient, the design extended beyond it: the emissions-related Mode~01 fuel PID returns a coarse, dashboard-damped percentage whose quantisation is inadequate for the $\pm 5\%$ target in SO1, so fuel level is instead read through the manufacturer-specific Toyota Mode~22 PID~21/29 channel at 0.5~L per ADC count. This access is manufacturer-specific rather than legislated and was validated on the partner Toyota vehicle, which also defines the portability boundary of the approach: extension to another make requires a per-vehicle adapter, and this is documented as a limitation of the design.

Two standards jointly fix the OBD sampling interval, and they act as opposing bounds. ISO~15765-4~\cite{iso15765_4} imposes a minimum inter-request interval, placing a ceiling on how fast the interface may be polled before the CAN transport saturates; SAE~J1349~\cite{sae_j1349}, by supplying the idle-combustion upper bound on fuel-loss rate, places a floor on how slowly it may be polled before the slosh envelope is under-sampled. The interval $\Delta s = 1.5$~s adopted in Sec.~5.4.3 is the value satisfying both constraints together with the median-window memory of Sec.~5.5.2, which makes the sampling rate an auditable derivation rather than a tuned parameter. ISO~7637-2~\cite{iso7637} produced a further consequence beyond power-supply selection: because transient-induced reset is an anticipated event on a vehicle supply line, buffered telemetry was required to survive reboot, which is why the offline buffer resides in SPIFFS flash rather than in RAM.

\paragraph{Data acquisition and signal processing.} The median--EMA--dead-band--direction-aware cascade in Sec.~5.5 is an engineering derivation from tank-slosh and sender-chatter behaviour; the standards observed here governed the structure of the record that the cascade produces. The FAIR principles translate, for a telemetry stream, into four schema requirements -- stable identifier, monotonic sequence number, dual timestamps, and channel provenance -- and these are what make offline-replayed records distinguishable from real-time records during evaluation. IETF RFC~5905 underwrites that schema by establishing the common time base: because the device clock is NTP-disciplined rather than free-running, device-side and backend-side timestamps are mutually comparable, drift is measurable as a reported quantity, and buffered records recover their correct ordering when replayed.

\paragraph{Communication and transmission.} A national regulatory instrument produced the decisive choice in this domain. A GSM-only module was the lower-cost option and adequate for a 5-second cadence, but NTC~MC No.~003-09-2025~\cite{ntc2025legacy_network_phaseout} provides for phased 2G and 3G decommissioning, and the LTE Cat-1 A7670E was therefore adopted so that LTE serves as the primary path with 2G retained as a transitional fallback. Two further decisions warrant explicit statements. First, observing ISO/IEC~27001~\cite{iso27001} through mandatory TLS means every wide-area upload incurs a negotiation cost, and the latency decomposition in Sec.~5.6.6 attributes a substantial share of the 23--24~s median LTE latency to it; confidentiality was prioritised over latency because telemetry carries driver location and behavioural data, and the trade-off is reported with its cause. Second, the OASIS MQTT specification was evaluated and deferred: with a single ingestion endpoint and no subscriber fan-out, its QoS~1 guarantees are met instead by HTTPS REST with the replay buffer, though the payload schema was kept broker-compatible so that a future migration requires no firmware rewrite.

\paragraph{Machine learning and validation.} ISO/IEC~29119~\cite{iso29119} shaped the test architecture rather than the models. Its requirement that test conditions be specified independently of the implementation produced three decisions: ground truth by controlled injection with known timing and magnitude, leave-one-trip-out cross-validation to prevent within-trip temporal leakage, and an out-of-distribution clean-data condition. The third is what gives the $\leq 10\%$ FPR claim in SO4 independent standing, since it bounds false alarms on data the detector never saw. IEEE~P7003~\cite{ieeep7003}, a draft project observed here as guidance, produced the explainability requirement: an anomaly flag is a decision about a driver, so alerts carry the contextual evidence that triggered them, and the contextual gate that suppresses alerts during refuelling and post-refuel quiescence is documented rather than implicit.

\paragraph{Software architecture and data management.} The value of observing ISO/IEC/IEEE~42010~\cite{iso42010} became evident when a constraint changed mid-development. Inference was originally intended to run on the ESP32, but available RAM proved insufficient for the Isolation Forest ensemble; because explicit interface contracts between layers had been maintained, inference was relocated to the backend while the acquisition, filtering, and transport layers remained untouched. The architecture description was therefore not documentation produced after the fact but the property that made a late structural change inexpensive. ISO/IEC~27018~\cite{iso27018}, together with GDPR principles, produced schema-level rather than policy-level decisions: driver identity is stored as an opaque reference, row-level security is enforced at the database layer so that an application defect cannot expose another operator's fleet, and dashboard views are role-scoped.

\paragraph{Basis of adherence.} The standards above were observed as design constraints, and each constraint is traceable to an implementation detail in Chapter~\ref{ch:method}. Adherence is therefore demonstrated through documented design decisions and through the measured results reported in Chapter~\ref{ch:result_discuss}, rather than through third-party conformity assessment, which was outside the scope of this study. Where a standard functions as a boundary rather than as a performance target, most notably ISO~26262~\cite{iso26262}, its contribution was to establish the limits of system authority.

\subsection{Observance of Research Ethics}

Ethical considerations guide the responsible deployment of the system, especially given the safety-critical and privacy-sensitive context of transportation fleets. Three major ethical domains were identified: data privacy, system safety, and responsible AI.

\subsubsection*{Data Privacy and Security}

The system adopts:

\begin{itemize}
    \item informed consent prior to any data collection,
    \item strict access control and secure data storage,
    \item anonymization and aggregation when feasible,
    \item transparency in data usage and access rights,
    \item compliance with national and international privacy regulations.
\end{itemize}

\subsubsection*{System Safety and Occupational Risk Management}

Given the operational risks of fleet environments, the design aligns with ISO~45001, ISO~26262, and ISO~39001 \cite{iso45001, iso26262, iso39001}. Safety-oriented measures include:

\begin{itemize}
    \item minimizing driver distraction,
    \item reducing false alarms through validation,
    \item providing adequate training for users,
    \item implementing clear protocols for anomaly response.
\end{itemize}

\subsubsection*{Responsible AI and Algorithmic Accountability}

Consistent with IEEE~P7003 \cite{ieeep7003}, the system incorporates:

\begin{itemize}
    \item explainability mechanisms for anomaly outputs,
    \item documentation of roles and responsibilities in alert handling,
    \item defined liability procedures for system errors,
    \item safeguards against algorithmic bias.
\end{itemize}

\section{Technical Design Considerations}  
\label{sec:techdesign}  

This section elaborates on the engineering decisions and standards that guided the development of the implemented IoT- and AI-driven fuel monitoring and anomaly detection system. Each component, software layer, algorithm, and communication protocol is grounded in established standards to ensure reliability, safety, interoperability, and security. The references indicate where each design decision is applied in the implementation procedures described in Chapter~\ref{ch:method}.  

\FloatBarrier
{\scriptsize
\setlength{\tabcolsep}{3pt}
\renewcommand{\arraystretch}{1.1}
\begin{longtable}{|p{0.14\textwidth}|p{0.19\textwidth}|p{0.17\textwidth}|p{0.27\textwidth}|p{0.17\textwidth}|}
\caption{Final Hardware and Subsystem Decision Table}
\label{tab:final_hardware_subsystem_decisions} \\
\hline
\textbf{Subsystem} & \textbf{Final choice} & \textbf{Alternative considered} & \textbf{Reason for final choice} & \textbf{Known limitation} \\
\hline
\endfirsthead

\hline
\textbf{Subsystem} & \textbf{Final choice} & \textbf{Alternative considered} & \textbf{Reason for final choice} & \textbf{Known limitation} \\
\hline
\endhead

Fuel acquisition & OBD-II Toyota Mode 22 PID 21/29 & External capacitive / float tank sensor & Non-invasive, no tank modification, ethically clearable, partner-acceptable & Manufacturer dashboard gauge used as field reference (no dipstick validation) \\
\hline
Controller & ESP32 (dual-core, 240 MHz, WiFi+BT\glsadd{BT}) & Arduino Uno / Raspberry Pi & Adequate GPIO + WiFi + dual core for IDF VSPI\glsadd{VSPI} separation; PlatformIO ecosystem; FreeRTOS scheduling & Limited RAM\glsadd{RAM} for on-device ML; ML inference moved to backend \\
\hline
HMI display & ILI9341 2.4" TFT + XPT2046 touch + dedicated buttons & OLED + scroll button only & Full route / fuel / alert UI\glsadd{UI} in-cab; suitable readability in daylight & Touch only used for non-critical actions; physical buttons used for safety-critical confirm \\
\hline
GPS & Embedded GPS (Neo-style) + mobile-app GPS fallback fused at backend & Dashboard-only / external assist only & Resilient to embedded-GPS cold start or building occlusion; tagged via \texttt{gps\_source} column & Urban-canyon drift; cold start of embedded module on first power-on \\
\hline
Wide-area comms - primary & LTE via A7670E modem & GSM-only baseline & Higher throughput, lower per-packet latency once attached & HTTPACTION + TLS renegotiation can push median latency to approximately 23--24 s on certain Smart cells \\
\hline
Wide-area comms - fallback & GSM 2G via same A7670E (with CGNAPN/CFUN+CSTT fallback) & LTE-only & Retained for legacy 2G-only coverage on some Batangas / Tagaytay routes & Lower throughput; not all carriers maintain 2G \\
\hline
Local-area / depot comms & LoRa SX1278 @ 433 MHz, SF7, BW 125 kHz, 17 dBm & WiFi-only at depot & Long range, low power, bidirectional fleet-yard polling without LTE cost & Low duty-cycle limits payload size; 125 kHz BW shared with other depot nodes \\
\hline
Hotspot fallback & WiFi (driver phone tether or depot AP) & No fallback & Recovers comms when LTE+GSM+LoRa are all down but a personal hotspot is available & Depends on driver compliance to enable hotspot \\
\hline
Offline buffer & SPIFFS append-only ring buffer, replayed with \texttt{channel\_used='offline\_replay'} & RAM-only buffer & Survives reboot; preserves \texttt{seq} + \texttt{event\_id} for backend deduplication & Flash wear if drain is starved for very long periods \\
\hline
Backend & Node/Express on Render + Supabase Postgres & Self-hosted Postgres on a VPS & Managed DB\glsadd{DB}, row-level security, free-tier acceptable for thesis-scale traffic & Render cold-start adds latency to the first POST after idleness \\
\hline
ML detector & Isolation Forest Model C (contamination=0.015, n\_estimators=300, max\_samples=256) & Threshold-only / IQR baseline; Hybrid Union; Hybrid Intersection & Only Isolation Forest variant that met $\geq$90\% precision and $\leq$10\% FPR under LOTO & Recall depends on injection realism; will be retrained as live anomaly corpus grows \\
\hline
Subsequence detector & Matrix Profile (window=12, z-thresh=2.0) & Drop entirely & Useful confirmatory layer for slow-leak / steady-siphon patterns that single-point IF misses & Less robust than Model C at row-level scoring; retained as supporting evidence only \\
\hline
\end{longtable}
}
\normalsize
\FloatBarrier

\subsection{Hardware Selection and Integration}  

The hardware configuration is designed to support the end-to-end workflow of the system (see Sec.~\ref{sec:implement}). The following standards and implementation points apply:

\begin{itemize}  
    \item \textbf{MCU:} ESP32/NodeMCU serves as the edge processing device for real-time acquisition and preprocessing of fuel, GPS, and vehicle telemetry data. It complies with IEC~61131-2 for embedded controllers \cite{iec61131} and ISO~7637 for automotive electrical transients \cite{iso7637}. In Sec.~\ref{sec:implement}, the MCU handles local preprocessing, timestamping, and alert-state handling before backend anomaly scoring.

    \item \textbf{GPS Module:} The UART-connected GPS module provides geolocation and speed data, adhering to global navigation standards \cite{gpsstd}. Its signals are synchronized with fuel readings during preprocessing in Sec.~\ref{sec:fuel_preprocessing}.

    \item \textbf{CAN/OBD-II Interface:} The CAN/OBD-II interface was selected as the acquisition layer because it satisfies the non-invasive constraint set by the partner operator. From this one port the prototype reads speed, RPM\glsadd{RPM}, engine load, throttle, MAF\glsadd{MAF}, coolant temperature, runtime, DTC\glsadd{DTC} status, and the Toyota Mode 22 PID 21/29 fuel-sender voltage encoded at 0.5 L per ADC\glsadd{ADC} count. Mode 22 access is manufacturer-specific and was validated on the partner Toyota vehicle.
\end{itemize}  

\subsection{Software Architecture}  

The modular software layers are designed according to ISO/IEC/IEEE~42010 software architecture standards \cite{iso42010}, supporting scalability, maintainability, and interoperability. Each layer’s function is mapped to its implementation in Sec.~\ref{sec:implement}:

\begin{itemize}  
    \item \textbf{Data Acquisition Layer:} Captures sensor and telemetry data, applies the implemented median + EMA + dead-band + direction-aware filter cascade, and prepares telemetry for backend anomaly scoring. Implementation details are in Secs.~\ref{sec:fuel_acquisition} and~\ref{sec:fuel_preprocessing}.  

    \item \textbf{Processing Layer:} Performs edge timestamping, trip-state handling, distance aggregation, and local alert-rule state preparation before cloud-side anomaly inference. Standards for real-time processing and task scheduling ensure predictable latency. Implementation is in Sec.~\ref{sec:fuel_preprocessing}.  

    \item \textbf{Cloud Integration Layer:} HTTPS REST endpoints receive JSON\glsadd{JSON} telemetry payloads from the prioritized communication stack. MQTT remains compatible as a future broker-based deployment option, while TLS/SSL ensures encrypted transmission (Sec.~\ref{sec:communication}).  

    \item \textbf{User Interface Layer:} Dashboards display real-time fuel levels, GPS locations, anomaly alerts, and historical trends. Web standards (HTML\glsadd{HTML}5, CSS3, JavaScript) and secure REST\glsadd{REST} APIs are employed, as described in Sec.~\ref{sec:backend}.  
\end{itemize}  

\subsection{Machine Learning and Algorithms}  

Machine learning components follow ISO/IEC~29119 software quality standards \cite{iso29119} and IEEE~P7003 AI accountability guidelines \cite{ieeep7003}:  

Static threshold fuel anomaly systems are insufficient in Philippine logistics because traffic congestion, idling, route delays, load variation, and stop-and-go driving can produce fuel-consumption changes that look abnormal under fixed rules. The proposed model uses contextual variables such as fuel rate, speed, RPM, engine load, throttle, idling, distance, and route behavior to distinguish suspicious fuel drops from normal but inefficient fuel use. This makes the system more adaptive than fixed threshold detection.

\begin{itemize}  
    \item \textbf{Models:} Isolation Forest was selected as the production detector because Model C (contamination = 0.015, \texttt{n\_estimators} = 300, \texttt{max\_samples} = 256, score threshold = the 0.985 quantile of the clean training scores) was the only configuration meeting both $\geq$90\% precision and $\leq$10\% FPR under leave-one-trip-out evaluation (precision = 92.9\%, recall = 92.9\%, FPR = 1.09\%; see Chapter~\ref{ch:result_discuss}). Matrix Profile (window = 12, z-threshold = 2.0) was retained as a supporting subsequence detector for slow-leak and siphon patterns that single-point scoring is less sensitive to.

    \item \textbf{Evaluation Metrics:} Confusion matrices, ROC curves, and precision–recall graphs are used for model validation, ensuring alignment with ISO/IEC 25010 software quality characteristics. Secs.~\ref{sec:evaluation} and~\ref{sec:evaluate} provide practical application in performance evaluation.

    \item \textbf{Alternative Models:} Autoencoders and LSTM\glsadd{LSTM} networks were considered but excluded due to edge processing constraints, as justified in Sec.~\ref{sec:alt_software_cost}.  
\end{itemize}  

\subsection{Communication and Security}  

Communication protocols and security measures are aligned with ISO/IEC 27018, ISO/IEC 27001, and GDPR standards to ensure data privacy, encryption, and secure access control:  

\begin{itemize}  
    \item \textbf{REST Telemetry:} HTTPS REST endpoints support direct cloud ingestion from LoRa gateway, LTE/GSM, WiFi, and offline-replay channels. MQTT with QoS\glsadd{QoS} 1 remains a future-compatible alternative if a broker-based deployment is adopted.

    \item \textbf{TLS/SSL Encryption:} Provides end-to-end encryption for transmitted data. Compliance with ISO/IEC~27001 ensures data integrity and confidentiality. Applied in Sec.~\ref{sec:communication}.  

    \item \textbf{Role-Based Access and Anonymization:} Enforces secure data access policies and anonymizes sensitive driver information, complying with ISO/IEC~27018 and GDPR principles. Applied in Sec.~\ref{sec:backend}.  

    \item \textbf{JSON Serialization:} Standardized data format for cloud compatibility, ensuring interoperability with cloud services and dashboards. Applied in Sec.~\ref{sec:communication}.  
\end{itemize}  

Wide-area connectivity uses a SIMCom A7670E LTE (Cat-1) module, chosen over a GSM-only baseline because it provides an LTE-primary, 2G-fallback hierarchy in one device. LTE is the primary path for its lowest latency and highest throughput, and because it is the network Philippine operators are expanding rather than retiring. GSM 2G on the same module is a deeper fallback for legacy-only coverage; consistent with NTC MC No. 003-09-2025 (phased 2G/3G decommissioning) it is treated as transitional, and because the A7670E is LTE-first rather than 2G/3G-only it remains compliant with that circular. LoRa provides depot, yard, and long-range telemetry; WiFi provides a depot or driver-tether path; and SPIFFS offline-buffer replay is the final fallback so no record is lost when all real-time channels are down \cite{ntc2025legacy_network_phaseout}. 

\section{Alternative Software Designs and Cost Considerations}
\label{sec:alt_software_cost}

During the planning and prototype development phase, several candidate software designs were evaluated for their suitability to the system’s requirements. Table~\ref{tab:alt_software} summarizes the strengths, potential risks, and fit relative to the system specifications. Matrix Profile (MP\glsadd{MP}), Isolation Forest (IF), and a combined Matrix Profile + Isolation Forest (MP + IF) configuration were implemented and tested to compare their effectiveness in detecting fuel-related anomalies. The comparison considered detection accuracy, computational efficiency, and deployment feasibility on the target hardware platform. The findings from these tests were used to determine the most appropriate approach for the final system implementation.

\begin{table}[h!]
\centering
\caption{Alternative Software Designs Evaluated for the Prototype}
\label{tab:alt_software}
\resizebox{\textwidth}{!}{%
\begin{tabular}{|l|p{4.5cm}|p{4.5cm}|c|}
\hline
\textbf{Candidate} & \textbf{Strengths} & \textbf{Risks / Trade-offs} & \textbf{Fit vs. Spec} \\
\hline
Isolation Forest (IF) & Fast, context-aware, handles multivariate data; Model C supports contextual fuel, speed, route, driver, and fuel-economy deviation features & Limited temporal context; sensitivity to contamination parameter may affect false alarms & Very High --- selected production detector \\
Matrix Profile (MP) & Shape-based time-series discord detection; useful for slow-leak and steady-siphon subsequence patterns as a confirmatory layer & Requires window-size tuning; primarily one-dimensional unless extended for multivariate data; less robust than Model C at row-level scoring & Moderate to High --- supporting subsequence detector \\
\hline
\end{tabular}%
}
\end{table}

A cost study was conducted after prototype testing to confirm that the system configuration was cost-effective. This study reviewed hardware, software, and development resources, comparing the chosen open-source and local execution setup against alternative configurations such as commercial analytics tools, cloud-based computing, or hardware-integrated telemetry systems. 

The total estimated budget for developing and implementing the prototype system is \textbf{PHP 10,446}, confirming that the selected configuration remains the optimal and sustainable choice for achieving the project’s objectives.


\section{Summary}
\label{sec:summary}

This chapter presented the key technical and ethical design considerations that shape the system. By grounding the project in ISO, IEEE, OASIS, GDPR, and FAIR standards, the system ensures reliable data acquisition, secure communication, validated machine learning, and responsible AI use.

These design elements establish the foundation for the procedures described in Chapter~\ref{ch:method}, bridging theoretical concepts from Chapter~\ref{ch:theorycon} with the practical implementation steps of the system.
